% ============================================================
% PAGE 4 — Problem Formulation & Proposed Architecture
% ============================================================

\section{Problem Formulation}
\label{sec:problem}

Let $\mathbf{I} \in \mathbb{R}^{3 \times H \times W}$ denote an
RGB input image of height $H$ and width $W$. The goal of object
detection is to infer a set of $K$ tuples
\begin{equation}
  \mathcal{O} = \bigl\{(b_k,\, c_k,\, s_k)\bigr\}_{k=1}^{K}
  \label{eq:detection_output}
\end{equation}
where $b_k = (x_k, y_k, w_k, h_k)$ specifies the axis-aligned
bounding box (center coordinates plus dimensions) of the $k$-th
detected instance, $c_k \in \{1, \ldots, C_{\text{cls}}\}$ is its
predicted class label from a vocabulary of $C_{\text{cls}}$ object
categories, and $s_k \in [0,1]$ is a confidence score indicating
the estimated probability that a genuine instance occupies the
predicted box.

A detection model $f_\theta$ with parameters $\theta$ implements this
mapping $f_\theta : \mathbf{I} \mapsto \mathcal{O}$ and is trained by
minimizing an expected composite loss
\begin{equation}
  \mathcal{L}(\theta) =
  \mathbb{E}_{(\mathbf{I},\mathcal{G})\sim\mathcal{D}}
  \bigl[\,
    \lambda_{\text{cls}}\,L_{\text{cls}}
    + \lambda_{\text{box}}\,L_{\text{box}}
    + \lambda_{\text{obj}}\,L_{\text{obj}}
  \,\bigr]
  \label{eq:total_loss}
\end{equation}
over the training distribution $\mathcal{D}$ of image--annotation
pairs $(\mathbf{I}, \mathcal{G})$, where $\mathcal{G}$ is the set of
ground-truth boxes and labels. The scalar weights $\lambda_{\text{cls}}$,
$\lambda_{\text{box}}$, and $\lambda_{\text{obj}}$ balance the three
loss terms; their values are determined by a grid search on the
validation split (Section~\ref{sec:results}).

\section{Proposed Architecture: FusionDet}
\label{sec:arch}

Figure~\ref{fig:arch_overview} provides a schematic of the complete
FusionDet pipeline, which consists of four sequential processing
stages: (A) the dual backbone, (B) the Adaptive Feature Fusion
module, (C) the bidirectional neck, and (D) the decoupled detection
head.

\begin{figure}[t]
\centering
\begin{tikzpicture}[
    box/.style   = {draw, rounded corners=2pt, minimum width=2.0cm,
                    minimum height=0.55cm, align=center, font=\small},
    arrow/.style = {-{Stealth[length=4pt]}, thick},
    scale=0.82, transform shape
  ]
  % Input
  \node[box, fill=gray!15] (inp) {Input Image};

  % CNN branch
  \node[box, fill=blue!15, below left=0.55cm and 0.5cm of inp]
    (cnn) {CSP-Darknet\\(CNN Branch)};
  % Transformer branch
  \node[box, fill=orange!15, below right=0.55cm and 0.5cm of inp]
    (vit) {Linear-Attn\\Transformer};

  % AFF
  \node[box, fill=green!20, below=1.1cm of inp]
    (aff) {Adaptive Feature\\Fusion (AFF)};

  % Neck
  \node[box, fill=purple!15, below=0.65cm of aff]
    (neck) {BiPANet Neck\\+ LCA};

  % Head
  \node[box, fill=red!15, below=0.65cm of neck]
    (head) {Decoupled\\Detection Head};

  % Output
  \node[box, fill=gray!15, below=0.55cm of head]
    (out) {Predictions $\mathcal{O}$};

  % Arrows
  \draw[arrow] (inp)  -- (cnn);
  \draw[arrow] (inp)  -- (vit);
  \draw[arrow] (cnn)  -- (aff);
  \draw[arrow] (vit)  -- (aff);
  \draw[arrow] (aff)  -- (neck);
  \draw[arrow] (neck) -- (head);
  \draw[arrow] (head) -- (out);
\end{tikzpicture}
\caption{High-level overview of the FusionDet architecture. Dual
backbones process the input in parallel; the AFF module merges their
outputs; the BiPANet neck builds a multi-scale pyramid; and the
decoupled head produces final predictions.}
\label{fig:arch_overview}
\end{figure}

\subsection{Dual-Backbone Design Rationale}
\label{sec:dual_backbone}

The two branches of FusionDet are designed to be architecturally
orthogonal:\ the CNN branch operates on local $3\times3$ kernels
with shift-equivariant inductive bias, while the transformer branch
uses linearized attention to aggregate global context from all
spatial positions simultaneously. Running these branches in parallel,
rather than composing one as a sub-module of the other, preserves the
maximum representational diversity.

Both branches accept the same $3\times H\times W$ input tensor and
independently produce four-level feature hierarchies at strides
$\{4, 8, 16, 32\}$. Let $\mathbf{F}^{\text{cnn}}_l$ and
$\mathbf{F}^{\text{vit}}_l$ denote the $l$-th level feature maps
from the CNN and transformer branches respectively, with
$l \in \{2,3,4,5\}$ following the standard FPN notation. The AFF
module is responsible for combining these pairs into a unified
representation $\mathbf{P}_l$ consumed by the neck.

\subsection{Adaptive Feature Fusion (AFF)}
\label{sec:aff}

Naive fusion strategies---element-wise addition, concatenation
followed by a $1\times1$ projection, or fixed weighted averages---
treat all channels and spatial locations identically. In practice,
some channels of the CNN output are more informative at small scales
(where convolutional edge detectors excel) while the transformer
output may dominate at large scales (where global context
disambiguates object identity). The AFF module learns this
trade-off through a pair of spatial--channel gating networks.

For each level $l$, the AFF module computes a gate tensor
$\mathbf{G}_l \in [0,1]^{C \times H_l \times W_l}$ from the
concatenated features:
\begin{equation}
  \mathbf{G}_l = \sigma\!\bigl(
    \mathcal{F}_{\text{gate}}\bigl(
      \text{Concat}(\mathbf{F}^{\text{cnn}}_l,\,
                     \mathbf{F}^{\text{vit}}_l)
    \bigr)
  \bigr)
  \label{eq:aff}
\end{equation}
where $\mathcal{F}_{\text{gate}}$ is a two-layer
$1\times1$--$3\times3$--$1\times1$ bottleneck with sigmoid output.
The fused feature map is then:
\begin{equation}
  \mathbf{P}_l =
    \mathbf{G}_l \odot \mathbf{F}^{\text{cnn}}_l
    \;+\;
    (1 - \mathbf{G}_l) \odot \mathbf{F}^{\text{vit}}_l
  \label{eq:aff_merge}
\end{equation}
This convex combination guarantees that the fused output lies in the
same feature space as each branch without the channel-dimension
doubling of concatenation.  Setting $\mathbf{G}_l \equiv 0.5$
recovers a simple average, confirming that the gated design strictly
generalizes simpler baselines.
